\documentclass[11pt,a4paper]{article}
\usepackage[utf8]{inputenc}
\usepackage[T1]{fontenc}
\usepackage[catalan]{babel}
\usepackage{lmodern}
\usepackage[sfdefault,lining]{FiraSans}
\usepackage[a4paper,top=2.3cm,bottom=2.2cm,left=1.8cm,right=1.8cm,headheight=1.3cm,headsep=0.4cm,footskip=1cm]{geometry}
\usepackage{amsmath,amssymb}
\usepackage{icomma}
\usepackage{graphicx}
\usepackage[table]{xcolor}
\usepackage{tikz}
\usepackage{fontawesome5}
\usepackage{ccicons}
\usepackage{fancyhdr}
\usepackage{lastpage}
\usepackage{tasks}
\usepackage{needspace}
\usepackage[most]{tcolorbox}
\usepackage{array}
\usepackage[hidelinks,pdftitle={Solucions de la fitxa de la unitat 2: Polinomis i fraccions algebraiques. 4t d'ESO Matemàtiques B},pdfsubject={Matemàtiques B 4t ESO},pdfkeywords={polinomis, operacions, Ruffini, valor numèric, teorema del residu, arrels, factorització, m.c.d., m.c.m., fraccions algebraiques}]{hyperref}

\definecolor{navy}{HTML}{1B2748}
\definecolor{accent}{HTML}{E8603A}
\definecolor{teal}{HTML}{1F9E90}
\definecolor{soft}{HTML}{EEF1F8}
\definecolor{muted}{HTML}{5D6685}

\setlength{\parindent}{0pt}
\setlength{\parskip}{3pt}
\settasks{label=\alph*), label-format=\bfseries\color{navy}, label-width=1.3em, item-indent=1.8em, column-sep=0.6em, after-item-skip=5pt, before-skip=3pt, after-skip=4pt}

% ---------- capçalera i peu ----------
\newcommand{\llicencia}{\href{https://creativecommons.org/licenses/by-nc-sa/4.0/deed.ca}{\ccbyncsa}\ \ Llicència Creative Commons Reconeixement-NoComercial-CompartirIgual 4.0 Internacional (CC BY-NC-SA 4.0)}
\pagestyle{fancy}
\fancyhf{}
\fancyhead[L]{\raisebox{-0.35\height}{\includegraphics[height=0.95cm]{logo-gva.png}}}
\fancyhead[C]{\small\color{muted}Solucions de la fitxa · Unitat 2: Polinomis i fraccions algebraiques · 4t ESO B}
\fancyhead[R]{\raisebox{-0.35\height}{\includegraphics[height=0.95cm]{logo-vermellar.png}}}
\fancyfoot[L]{\scriptsize\color{muted}\llicencia}
\fancyfoot[R]{\small\color{muted}\thepage}
\renewcommand{\headrulewidth}{0.4pt}
\renewcommand{\headrule}{\color{muted!40}\hrule height \headrulewidth}
\fancypagestyle{primera}{\fancyhf{}\renewcommand{\headrulewidth}{0pt}
  \fancyfoot[L]{\scriptsize\color{muted}\llicencia}
  \fancyfoot[R]{\small\color{muted}\thepage}}

% ---------- nivells i exercicis ----------
\newcommand{\estrelles}[1]{{\color{accent}\ifcase#1\or\faStar\faStar[regular]\faStar[regular]\or\faStar\faStar\faStar[regular]\or\faStar\faStar\faStar\fi}}
\newcounter{exer}
\newcommand{\ex}[1]{\par\needspace{4\baselineskip}\medskip\refstepcounter{exer}\noindent{\color{navy}\bfseries\theexer.}\ \ {\small\estrelles{#1}}\ \ }

\newcommand{\seccio}[2]{\par\needspace{16\baselineskip}\bigskip
  \begin{tcolorbox}[colback=navy,colframe=navy,coltext=white,arc=2mm,boxrule=0pt,left=3mm,right=3mm,top=1.2mm,bottom=1.2mm]
  {\large\bfseries #1\quad #2}
  \end{tcolorbox}\vspace{-2pt}}

\newcommand{\exf}[1]{\par\medskip\refstepcounter{exer}\noindent{\color{navy}\bfseries\theexer.}\ \ {\small\estrelles{#1}}\ \ }
\newcolumntype{C}[1]{>{\centering\arraybackslash}p{#1}}
% quadre de recordatori (només a la fitxa d'exercicis)
\newcommand{\recorda}[1]{\par\smallskip\begin{tcolorbox}[colback=accent!7,colframe=accent!55,arc=2mm,boxrule=0.5pt,left=2.5mm,right=2.5mm,top=1mm,bottom=1mm]\small{\color{accent}\faLightbulb[regular]}\ \textbf{Recorda:}\ #1\end{tcolorbox}}
\newcommand{\res}[1]{\boxed{#1}}\newcommand{\rest}[1]{\fbox{#1}}

\def\fmneed{11}% que el títol d'una secció no quede sol al final del full
% =====================================================================
% Capçalera FentMates de les fitxes: un color per a cada curs.
% Els colors són els de la pàgina inicial (LEVELS a assets/js/materials.js);
% el logo ve de _eines/logo/pdf/ (python3 _eines/logo/genera.py).
%
% Ús, al final del preàmbul de la unitat:
%   % =====================================================================
% Capçalera FentMates de les fitxes: un color per a cada curs.
% Els colors són els de la pàgina inicial (LEVELS a assets/js/materials.js);
% el logo ve de _eines/logo/pdf/ (python3 _eines/logo/genera.py).
%
% Ús, al final del preàmbul de la unitat:
%   % =====================================================================
% Capçalera FentMates de les fitxes: un color per a cada curs.
% Els colors són els de la pàgina inicial (LEVELS a assets/js/materials.js);
% el logo ve de _eines/logo/pdf/ (python3 _eines/logo/genera.py).
%
% Ús, al final del preàmbul de la unitat:
%   \input{../comu/fentmates.tex}
%   \fmcurs{1eso}
% i, al començament del document:
%   \capcalera{Matemàtiques}{1r ESO}{Fitxa d'exercicis · Unitat 2}{Divisibilitat}
% Opcional abans de \input: \def\fmneed{14} (espai mínim abans de cada secció).
% =====================================================================
\makeatletter
\def\fm@logodir{../../logo/pdf}
\providecommand{\fmneed}{10}
\def\fm@cursnom{}

% clau del curs, color i color del text damunt del color
\newcommand{\fm@curs}[3]{\@namedef{fm@color@#1}{#2}\@namedef{fm@text@#1}{#3}}
\fm@curs{1eso}{E6B84A}{navy}
\fm@curs{2eso}{3B82F6}{white}
\fm@curs{3eso}{FF7A50}{white}
\fm@curs{4eso-a}{A78BFA}{white}
\fm@curs{4eso-b}{6D28D9}{white}
\fm@curs{1batx-mi}{9B1C1C}{white}
\fm@curs{1batx-ccss}{EF6B6B}{white}
\fm@curs{2batx-mii}{1E6B47}{white}
\fm@curs{2batx-ccss}{2BB3A3}{white}

\newcommand{\fmcurs}[1]{%
  \@ifundefined{fm@color@#1}{\PackageError{fentmates}{Curs desconegut: #1}{}}{}%
  \edef\fm@tmp{\noexpand\definecolor{curs}{HTML}{\@nameuse{fm@color@#1}}%
               \noexpand\colorlet{curstext}{\@nameuse{fm@text@#1}}}\fm@tmp
  \colorlet{cursfosc}{curs!72!black}%
  \def\fm@logo{\fm@logodir/vertical-#1.pdf}%
  \def\fm@marca{\fm@logodir/marca-#1.pdf}}

% ---------- capçalera de la primera pàgina ----------
% #1 matèria (petit), #2 curs, #3 tipus de document i unitat, #4 títol
\newcommand{\capcalera}[4]{%
  \gdef\fm@cursnom{#2}%
  \par\noindent
  \begin{tcolorbox}[colback=curs,colframe=curs,boxrule=0pt,arc=2.2mm,width=3.5cm,height=2.15cm,
      valign=center,left=3mm,right=2mm,top=1mm,bottom=1mm,nobeforeafter,box align=center]
    \color{curstext}{\footnotesize #1}\par\vspace{1pt}{\fontsize{19}{21}\selectfont\bfseries #2\par}%
  \end{tcolorbox}\hspace{5mm}%
  \begin{minipage}[c]{\dimexpr\linewidth-3.5cm-5mm-2.9cm\relax}
    {\color{cursfosc}\small\bfseries\MakeUppercase{#3}}\par\vspace{2pt}
    {\color{navy}\fontsize{20}{23}\selectfont\bfseries #4\par}
  \end{minipage}\hfill
  \begin{minipage}[c]{2.6cm}\raggedleft\includegraphics[width=2.6cm]{\fm@logo}\end{minipage}\par
  \vspace{2pt}}

% ---------- barres de secció del color del curs ----------
\renewcommand{\seccio}[2]{\par\needspace{\fmneed\baselineskip}\bigskip
  \begin{tcolorbox}[colback=curs,colframe=curs,coltext=curstext,arc=2mm,boxrule=0pt,left=3mm,right=3mm,top=1.2mm,bottom=1.2mm]
  {\large\bfseries #1\quad #2}
  \end{tcolorbox}\vspace{-2pt}}

% ---------- estrelles de dificultat del color del curs (a Batxillerat no se'n posen) ----------
\@ifundefined{estrelles}{}{%
  \renewcommand{\estrelles}[1]{{\color{curs}\ifcase#1\or\faStar\faStar[regular]\faStar[regular]\or\faStar\faStar\faStar[regular]\or\faStar\faStar\faStar\fi}}}

% ---------- capçalera i peu de les pàgines ----------
% primera pàgina: logos institucionals al peu
\fancypagestyle{primera}{\fancyhf{}\renewcommand{\headrulewidth}{0pt}
  \fancyfoot[L]{\raisebox{-0.32\height}{\includegraphics[height=0.78cm]{logo-gva.png}}\hspace{5mm}%
    \raisebox{-0.32\height}{\includegraphics[height=0.78cm]{logo-vermellar.png}}}
  \fancyfoot[R]{\scriptsize\color{muted}\href{https://creativecommons.org/licenses/by-nc-sa/4.0/deed.ca}{\ccbyncsa}\ \ CC BY-NC-SA 4.0 · fentmates.org\qquad\small\thepage}}
% la resta: el símbol a l'esquerra i el curs a la dreta, amb el títol al centre com abans
\fancyhead[L]{\raisebox{-0.3\height}{\includegraphics[height=0.62cm]{\fm@marca}}}
\fancyhead[R]{\small\bfseries\color{cursfosc}\fm@cursnom}
\renewcommand{\headrule}{\color{curs!45}\hrule height \headrulewidth}
\makeatother

%   \fmcurs{1eso}
% i, al començament del document:
%   \capcalera{Matemàtiques}{1r ESO}{Fitxa d'exercicis · Unitat 2}{Divisibilitat}
% Opcional abans de \input: \def\fmneed{14} (espai mínim abans de cada secció).
% =====================================================================
\makeatletter
\def\fm@logodir{../../logo/pdf}
\providecommand{\fmneed}{10}
\def\fm@cursnom{}

% clau del curs, color i color del text damunt del color
\newcommand{\fm@curs}[3]{\@namedef{fm@color@#1}{#2}\@namedef{fm@text@#1}{#3}}
\fm@curs{1eso}{E6B84A}{navy}
\fm@curs{2eso}{3B82F6}{white}
\fm@curs{3eso}{FF7A50}{white}
\fm@curs{4eso-a}{A78BFA}{white}
\fm@curs{4eso-b}{6D28D9}{white}
\fm@curs{1batx-mi}{9B1C1C}{white}
\fm@curs{1batx-ccss}{EF6B6B}{white}
\fm@curs{2batx-mii}{1E6B47}{white}
\fm@curs{2batx-ccss}{2BB3A3}{white}

\newcommand{\fmcurs}[1]{%
  \@ifundefined{fm@color@#1}{\PackageError{fentmates}{Curs desconegut: #1}{}}{}%
  \edef\fm@tmp{\noexpand\definecolor{curs}{HTML}{\@nameuse{fm@color@#1}}%
               \noexpand\colorlet{curstext}{\@nameuse{fm@text@#1}}}\fm@tmp
  \colorlet{cursfosc}{curs!72!black}%
  \def\fm@logo{\fm@logodir/vertical-#1.pdf}%
  \def\fm@marca{\fm@logodir/marca-#1.pdf}}

% ---------- capçalera de la primera pàgina ----------
% #1 matèria (petit), #2 curs, #3 tipus de document i unitat, #4 títol
\newcommand{\capcalera}[4]{%
  \gdef\fm@cursnom{#2}%
  \par\noindent
  \begin{tcolorbox}[colback=curs,colframe=curs,boxrule=0pt,arc=2.2mm,width=3.5cm,height=2.15cm,
      valign=center,left=3mm,right=2mm,top=1mm,bottom=1mm,nobeforeafter,box align=center]
    \color{curstext}{\footnotesize #1}\par\vspace{1pt}{\fontsize{19}{21}\selectfont\bfseries #2\par}%
  \end{tcolorbox}\hspace{5mm}%
  \begin{minipage}[c]{\dimexpr\linewidth-3.5cm-5mm-2.9cm\relax}
    {\color{cursfosc}\small\bfseries\MakeUppercase{#3}}\par\vspace{2pt}
    {\color{navy}\fontsize{20}{23}\selectfont\bfseries #4\par}
  \end{minipage}\hfill
  \begin{minipage}[c]{2.6cm}\raggedleft\includegraphics[width=2.6cm]{\fm@logo}\end{minipage}\par
  \vspace{2pt}}

% ---------- barres de secció del color del curs ----------
\renewcommand{\seccio}[2]{\par\needspace{\fmneed\baselineskip}\bigskip
  \begin{tcolorbox}[colback=curs,colframe=curs,coltext=curstext,arc=2mm,boxrule=0pt,left=3mm,right=3mm,top=1.2mm,bottom=1.2mm]
  {\large\bfseries #1\quad #2}
  \end{tcolorbox}\vspace{-2pt}}

% ---------- estrelles de dificultat del color del curs (a Batxillerat no se'n posen) ----------
\@ifundefined{estrelles}{}{%
  \renewcommand{\estrelles}[1]{{\color{curs}\ifcase#1\or\faStar\faStar[regular]\faStar[regular]\or\faStar\faStar\faStar[regular]\or\faStar\faStar\faStar\fi}}}

% ---------- capçalera i peu de les pàgines ----------
% primera pàgina: logos institucionals al peu
\fancypagestyle{primera}{\fancyhf{}\renewcommand{\headrulewidth}{0pt}
  \fancyfoot[L]{\raisebox{-0.32\height}{\includegraphics[height=0.78cm]{logo-gva.png}}\hspace{5mm}%
    \raisebox{-0.32\height}{\includegraphics[height=0.78cm]{logo-vermellar.png}}}
  \fancyfoot[R]{\scriptsize\color{muted}\href{https://creativecommons.org/licenses/by-nc-sa/4.0/deed.ca}{\ccbyncsa}\ \ CC BY-NC-SA 4.0 · fentmates.org\qquad\small\thepage}}
% la resta: el símbol a l'esquerra i el curs a la dreta, amb el títol al centre com abans
\fancyhead[L]{\raisebox{-0.3\height}{\includegraphics[height=0.62cm]{\fm@marca}}}
\fancyhead[R]{\small\bfseries\color{cursfosc}\fm@cursnom}
\renewcommand{\headrule}{\color{curs!45}\hrule height \headrulewidth}
\makeatother

%   \fmcurs{1eso}
% i, al començament del document:
%   \capcalera{Matemàtiques}{1r ESO}{Fitxa d'exercicis · Unitat 2}{Divisibilitat}
% Opcional abans de \input: \def\fmneed{14} (espai mínim abans de cada secció).
% =====================================================================
\makeatletter
\def\fm@logodir{../../logo/pdf}
\providecommand{\fmneed}{10}
\def\fm@cursnom{}

% clau del curs, color i color del text damunt del color
\newcommand{\fm@curs}[3]{\@namedef{fm@color@#1}{#2}\@namedef{fm@text@#1}{#3}}
\fm@curs{1eso}{E6B84A}{navy}
\fm@curs{2eso}{3B82F6}{white}
\fm@curs{3eso}{FF7A50}{white}
\fm@curs{4eso-a}{A78BFA}{white}
\fm@curs{4eso-b}{6D28D9}{white}
\fm@curs{1batx-mi}{9B1C1C}{white}
\fm@curs{1batx-ccss}{EF6B6B}{white}
\fm@curs{2batx-mii}{1E6B47}{white}
\fm@curs{2batx-ccss}{2BB3A3}{white}

\newcommand{\fmcurs}[1]{%
  \@ifundefined{fm@color@#1}{\PackageError{fentmates}{Curs desconegut: #1}{}}{}%
  \edef\fm@tmp{\noexpand\definecolor{curs}{HTML}{\@nameuse{fm@color@#1}}%
               \noexpand\colorlet{curstext}{\@nameuse{fm@text@#1}}}\fm@tmp
  \colorlet{cursfosc}{curs!72!black}%
  \def\fm@logo{\fm@logodir/vertical-#1.pdf}%
  \def\fm@marca{\fm@logodir/marca-#1.pdf}}

% ---------- capçalera de la primera pàgina ----------
% #1 matèria (petit), #2 curs, #3 tipus de document i unitat, #4 títol
\newcommand{\capcalera}[4]{%
  \gdef\fm@cursnom{#2}%
  \par\noindent
  \begin{tcolorbox}[colback=curs,colframe=curs,boxrule=0pt,arc=2.2mm,width=3.5cm,height=2.15cm,
      valign=center,left=3mm,right=2mm,top=1mm,bottom=1mm,nobeforeafter,box align=center]
    \color{curstext}{\footnotesize #1}\par\vspace{1pt}{\fontsize{19}{21}\selectfont\bfseries #2\par}%
  \end{tcolorbox}\hspace{5mm}%
  \begin{minipage}[c]{\dimexpr\linewidth-3.5cm-5mm-2.9cm\relax}
    {\color{cursfosc}\small\bfseries\MakeUppercase{#3}}\par\vspace{2pt}
    {\color{navy}\fontsize{20}{23}\selectfont\bfseries #4\par}
  \end{minipage}\hfill
  \begin{minipage}[c]{2.6cm}\raggedleft\includegraphics[width=2.6cm]{\fm@logo}\end{minipage}\par
  \vspace{2pt}}

% ---------- barres de secció del color del curs ----------
\renewcommand{\seccio}[2]{\par\needspace{\fmneed\baselineskip}\bigskip
  \begin{tcolorbox}[colback=curs,colframe=curs,coltext=curstext,arc=2mm,boxrule=0pt,left=3mm,right=3mm,top=1.2mm,bottom=1.2mm]
  {\large\bfseries #1\quad #2}
  \end{tcolorbox}\vspace{-2pt}}

% ---------- estrelles de dificultat del color del curs (a Batxillerat no se'n posen) ----------
\@ifundefined{estrelles}{}{%
  \renewcommand{\estrelles}[1]{{\color{curs}\ifcase#1\or\faStar\faStar[regular]\faStar[regular]\or\faStar\faStar\faStar[regular]\or\faStar\faStar\faStar\fi}}}

% ---------- capçalera i peu de les pàgines ----------
% primera pàgina: logos institucionals al peu
\fancypagestyle{primera}{\fancyhf{}\renewcommand{\headrulewidth}{0pt}
  \fancyfoot[L]{\raisebox{-0.32\height}{\includegraphics[height=0.78cm]{logo-gva.png}}\hspace{5mm}%
    \raisebox{-0.32\height}{\includegraphics[height=0.78cm]{logo-vermellar.png}}}
  \fancyfoot[R]{\scriptsize\color{muted}\href{https://creativecommons.org/licenses/by-nc-sa/4.0/deed.ca}{\ccbyncsa}\ \ CC BY-NC-SA 4.0 · fentmates.org\qquad\small\thepage}}
% la resta: el símbol a l'esquerra i el curs a la dreta, amb el títol al centre com abans
\fancyhead[L]{\raisebox{-0.3\height}{\includegraphics[height=0.62cm]{\fm@marca}}}
\fancyhead[R]{\small\bfseries\color{cursfosc}\fm@cursnom}
\renewcommand{\headrule}{\color{curs!45}\hrule height \headrulewidth}
\makeatother

\fmcurs{4eso-b}

\begin{document}
\thispagestyle{primera}

\capcalera{Matemàtiques B}{4t ESO}{Solucions de la fitxa · Unitat 2}{\raggedright Polinomis i fraccions algebraiques}
\vspace{8pt}

\begin{tcolorbox}[colback=soft,colframe=soft,arc=2mm,left=3mm,right=3mm,top=1.5mm,bottom=1.5mm]
\small Intenta resoldre cada exercici abans de mirar-ne la solució. Pots comprovar les divisions, Ruffini, les arrels, les factoritzacions i les fraccions amb les eines interactives dels apunts de la unitat.
\end{tcolorbox}

% =====================================================
\seccio{1}{Polinomis}

\exf{1} \begin{tasks}(1)
  \task $\res{2x^{4}+3x^{2}-x-5}$: grau 4, coeficient principal 2, terme independent $-5$. No és complet (falta $x^3$).
  \task $\res{-x^{3}+7}$: grau 3, coeficient principal $-1$, terme independent 7. No és complet.
  \task $\res{x^{5}+x^{3}-2x^{2}-8x+1}$: grau 5, coeficient principal 1, terme independent 1. No és complet (falta $x^4$).
  \task $\res{6x^{2}-4x-\frac{1}{2}}$: grau 2, coeficient principal 6, terme independent $-\frac12$. És complet.
\end{tasks}

\ex{1} \begin{tasks}(2)
  \task $-3$, \ $x^4$, \ $\res{4}$
  \task $\frac25$, \ $xy^2$, \ $\res{3}$
  \task $8$, \ cap, \ $\res{0}$
  \task $-1$, \ $a^3b^2$, \ $\res{5}$
\end{tasks}

\ex{1} \begin{tasks}(3)
  \task $\res{2x^{2}-3x+3}$
  \task $\res{5x^{3}+x^{2}-3x}$
  \task $\res{x^{2}-6x}$
\end{tasks}


% =====================================================
\seccio{2}{Operacions amb polinomis}

\exf{1} \begin{tasks}(2)
  \task $\res{2x^{3}+3x+3}$
  \task $\res{2x^{3}-2x^{2}-3x+5}$
  \task $\res{x^{3}+x^{2}-7x+2}$
  \task $\res{2x^{4}-5x^{3}+2x^{2}+4x-8}$
  \task $\res{x^{3}-2x^{2}+7x+2}$
\end{tasks}

\ex{1} \begin{tasks}(3)
  \task $\res{x^{2}+10x+25}$
  \task $\res{4x^{2}-12x+9}$
  \task $\res{9x^{2}-1}$
  \task $\res{x^{4}-4x^{2}+4}$
  \task $\res{\frac{1}{4}x^{2}+4x+16}$
  \task $\res{-x^{2}+25}$
\end{tasks}

\ex{2} \begin{tasks}(4)
  \task $\res{(x+3)^2}$
  \task $\res{(2x-5)^2}$
  \task $\res{(x+7)(x-7)}$
  \task $\res{(3x+1)(3x-1)}$
\end{tasks}

\ex{2} \begin{tasks}(3)
  \task $\res{8x}$
  \task $\res{-x^{2}-6x+12}$
  \task $\res{2x^{2}-4x+3}$
\end{tasks}

\ex{2} \begin{tasks}(1)
  \task $c(x)=\res{x-2},\quad r(x)=\res{2x+1}$
  \task $c(x)=\res{3x^{2}-2},\quad r(x)=\res{3x+2}$
  \task $c(x)=\res{2x^{2}+3x+1},\quad r(x)=\res{-7x+1}$
  \task $c(x)=\res{x^{2}-1},\quad r(x)=\res{0}$
\end{tasks}

\ex{1} \begin{tasks}(1)
  \task $c(x)=\res{x^{2}-2x-3},\quad r=\res{0}$
  \task $c(x)=\res{2x^{3}-2x^{2}-x+2},\quad r=\res{-7}$
  \task $c(x)=\res{x^{4}+2x^{3}+4x^{2}+8x+16},\quad r=\res{0}$
  \task $c(x)=\res{3x^{2}-4x+7},\quad r=\res{-14}$
\end{tasks}

\ex{2} $(x+a)(2x-3)=2x^2+(2a-3)x-3a$. Igualant coeficients: $-3a=-6\Rightarrow\res{a=2}$ \ i \ $b=2a-3\Rightarrow\res{b=1}$.


% =====================================================
\seccio{3}{Valor numèric i teorema del residu}

\exf{1} \begin{tasks}(5)
  \task $\res{-1}$
  \task $\res{-1}$
  \task $\res{-7}$
  \task $\res{-1}$
  \task $\res{-\frac{5}{8}}$
\end{tasks}

\ex{2} \begin{tasks}(2)
  \task $P(1)=1-3+2=\res{0}$
  \task $P(-1)=-1-2-3=\res{-6}$
  \task $P(2)=16-2+5=\res{19}$
  \task $P(-1)=1-1=\res{0}$
\end{tasks}

\ex{2} Ruffini amb 2: coeficients $1,\ 0,\ -5,\ 3$ $\to$ $1,\ 2,\ -1,\ 1$; el residu és $\res{1}$. I $P(2)=8-10+3=1$. Coincideixen.

\ex{2} \begin{tasks}(1)
  \task $P(3)=27+9k-15+6=0\Rightarrow\res{k=-2}$
  \task $P(-2)=-16+2k+4=8\Rightarrow\res{k=10}$
  \task $P(-2)=16+4k-4=0\Rightarrow\res{k=-3}$
\end{tasks}

\ex{3} $P(2)=0$ i $P(-3)=0$: $\begin{cases}4a+2b=-2\\ 9a-3b=33\end{cases}\Rightarrow\res{a=2,\ b=-5}$.


% =====================================================
\seccio{4}{Arrels d'un polinomi}

\exf{1} \begin{tasks}(4)
  \task $P(1)=0$: $\res{\text{sí}}$
  \task $P(2)=0$: $\res{\text{sí}}$
  \task $P(-1)=12$: $\res{\text{no}}$
  \task $P(-3)=0$: $\res{\text{sí}}$
\end{tasks}

\ex{2} \begin{tasks}(2)
  \task $\res{x=1,\ x=2,\ x=3}$
  \task $\res{x=1,\ x=-1,\ x=-2}$
  \task $\res{x=\pm1,\ x=\pm3}$
  \task $\res{x=2,\ x=-2,\ x=3}$
\end{tasks}

\ex{2} Divisors de 12: $\pm1,\ \pm2,\ \pm3,\ \pm4,\ \pm6,\ \pm12$. Només $P(2)=8+4-12=0$: l'única arrel és $\res{x=2}$ (el quocient $x^2+3x+6$ no té arrels reals).

\ex{2} $(x-1)(x+2)(x-3)=\res{x^{3}-2x^{2}-5x+6}$; \ $(x-4)^2=\res{x^{2}-8x+16}$ (o qualsevol múltiple).


% =====================================================
\seccio{5}{Factorització de polinomis}

\exf{1} \begin{tasks}(3)
  \task $\res{3x(x-4)}$
  \task $\res{x(x+3)(x-3)}$
  \task $\res{2x(x+2)^2}$
  \task $\res{x^2(x+1)(x-1)}$
  \task $\res{5(x+1)(x-1)}$
\end{tasks}

\ex{2} \begin{tasks}(2)
  \task $\res{(x-2)(x-3)}$
  \task $\res{(x+4)(x-3)}$
  \task $\res{2\left(x+\tfrac12\right)(x-1)=(2x+1)(x-1)}$
  \task Irreductible ($\Delta=-4<0$).
  \task $\res{3(x+3)(x-1)}$
\end{tasks}

\ex{2} \begin{tasks}(2)
  \task $\res{(x-1)(x+2)(x-3)}$
  \task $\res{(x-1)^2(x+3)}$
  \task $\res{(x-1)(x+1)(x+2)(x-3)}$
  \task $\res{(x+1)^2(x-2)}$
  \task $\res{(x-1)(x+1)(x-2)(x+2)}$
  \task $\res{(x-1)^2(x+1)}$
\end{tasks}

\ex{3} \begin{tasks}(2)
  \task $\res{x^2(x-3)(x+1)}$
  \task $\res{x(x-1)(2x-1)(x+2)}$
  \task $\res{x(x-1)(x+1)(x^2+1)}$
  \task $\res{(x-1)(x+1)(x^2+x+1)}$
\end{tasks}

\ex{1} $x=0$ ($\res{\text{doble}}$), $x=3$ ($\res{\text{triple}}$), $x=-1$ ($\res{\text{simple}}$). Grau $2+3+1=\res{6}$.


% =====================================================
\seccio{6}{Divisibilitat. M.c.d. i m.c.m.}

\exf{1} \begin{tasks}(2)
  \task $1-1=0$: $\res{\text{sí}}$
  \task $4+4=8\neq0$: $\res{\text{no}}$
  \task $1-1=0$: $\res{\text{sí}}$
  \task $9-9=0$: $\res{\text{sí}}$
\end{tasks}

\ex{2} \begin{tasks}(1)
  \task $(x+2)(x-2)$, \ $(x+2)^2$: \ m.c.d. $=\res{x+2}$, \ m.c.m. $=\res{(x+2)^2(x-2)}$
  \task $x(x-1)$, \ $(x+1)(x-1)$: \ m.c.d. $=\res{x-1}$, \ m.c.m. $=\res{x(x+1)(x-1)}$
  \task $x(x+1)(x-1)$, \ $(x-1)^2$: \ m.c.d. $=\res{x-1}$, \ m.c.m. $=\res{x(x+1)(x-1)^2}$
  \task $2(x+2)$, \ $x(x+2)$, \ $3x$: \ m.c.d. $=\res{1}$, \ m.c.m. $=\res{6x(x+2)}$
\end{tasks}

\ex{3} $P(x)=(x-1)(x+1)(x-2)$, \ $Q(x)=(x-1)(x-2)$. \ m.c.d. $=\res{(x-1)(x-2)}$, \ m.c.m. $=\res{(x-1)(x+1)(x-2)}$.


% =====================================================
\seccio{7}{Fraccions algebraiques}

\exf{1} \begin{tasks}(4)
  \task $\res{x=2}$
  \task $\res{x=\pm3}$
  \task Sempre en tenen.
  \task $\res{x=0,\ x=4}$
\end{tasks}

\ex{2} \begin{tasks}(3)
  \task $\res{\dfrac{x-3}{x}}$
  \task $\res{\dfrac{x-2}{x+2}}$
  \task $\res{\dfrac{x-1}{2}}$
  \task $\res{x-1}$
  \task $\res{\dfrac{x-2}{x+2}}$
\end{tasks}

\ex{2} \begin{tasks}(1)
  \task $\res{\text{Sí}}$: $(x+1)(x^2-1)=(x-1)(x^2+2x+1)$ (o bé, simplificant la segona: $\frac{(x+1)^2}{(x+1)(x-1)}=\frac{x+1}{x-1}$).
  \task $\res{\text{No}}$: $2(x^2+1)\neq x(2x+2)$.
\end{tasks}

\ex{2} \begin{tasks}(2)
  \task $\res{\dfrac{2x+3}{x^2}}$
  \task $\res{\dfrac{2x}{x^2-1}}$
  \task $\res{\dfrac{x^2+4}{x^2-4}}$
  \task $\res{\dfrac{4-x}{x^2-1}}$
  \task $\res{\dfrac{2x+3}{x(x+1)}}$
\end{tasks}

\ex{2} \begin{tasks}(2)
  \task $\res{2x(x-2)}$
  \task $\res{\dfrac1x}$
  \task $\res{(x-1)(x-3)}$
  \task $\res{\dfrac{1}{2x(x-1)}}$
\end{tasks}

\ex{3} \begin{tasks}(2)
  \task $\res{\dfrac{1}{x-1}}$
  \task $\res{x+1}$
  \task $\res{1}$
  \task $\res{\dfrac{8x}{x^2-4}}$
\end{tasks}


% =====================================================
\seccio{8}{Problemes}

\exf{2} Altura: $(x^2+5x+6):(x+3)=\res{x+2\ \text{cm}}$. Perímetre: $2(x+3)+2(x+2)=\res{4x+10\ \text{cm}}$.

\ex{2} $\dfrac{x^3+6x^2+8x}{x(x+2)}=\dfrac{x(x+2)(x+4)}{x(x+2)}=\res{x+4\ \text{cm}}$.

\ex{1} $(x^3-x+5)-(3x^2-2x+1)=\res{x^{3}-3x^{2}+x+4}$.


\end{document}
